\section{Relational norms and the Cohen value calculation}
\label{annex:p06-cardinal-characteristics-separations}

\begin{annexitem}{annex:p06.ccsep.morphism-duality}{thm:p06-cc-morphism-duality}{Transporting witnesses in both directions}
\begin{proof}
If \(Z\subseteq A_+\) witnesses \(\|\mathbf A\|\), then for every
\(b\in B_-\) choose \(z\in Z\) with \(A(\xi_-(b),z)\).
The morphism implication gives \(B(b,\xi_+(z))\), so
\(\xi_+[Z]\) witnesses \(\|\mathbf B\|\) and has size at most \(|Z|\).
For duality take the contrapositive of the same implication:
\[
 \neg B(b,\xi_+(a))\ \Longrightarrow\ \neg A(\xi_-(b),a).
\]
The negative component for \(\mathbf B^\perp\to \mathbf A^\perp\) is
\(\xi_+:A_+\to B_+\), and its positive component is
\(\xi_-:B_-\to A_-\). The displayed contrapositive is exactly
its required implication. Applying the first witness argument to this
dual morphism gives the second norm inequality.

For the interval construction, use
\(\mathbf W=(\omega^\omega,\omega^\omega,\leq^*)\) and the splitting
datum \(\mathbf S=(\mathcal P_\infty(\omega),\mathcal P(\omega),
\text{is split by})\). This letter names only the local datum in this proof.
The map \(\xi_-(Y)\) sends \(k\) to the next member
of \(Y\) after \(k\), while \(\xi_+(f)\) is the union of even
intervals with endpoints \(a_0=0\),
\(a_{n+1}=1+\max\{a_n,f(a_n)\}\).
The already supplied interval proof establishes
\(\xi_-(Y)\leq^*f\Rightarrow Y\) is split by \(\xi_+(f)\).
This is a morphism \(\mathbf W\to\mathbf S\). The norm and dual-norm readings
in the definition give the two asserted cardinal inequalities.
\end{proof}
\end{annexitem}

\begin{annexitem}{annex:p06.ccsep.sequential-norm}{thm:p06-cc-sequential-norm}{The full product and dual-minimum argument}
\begin{proof}
Let \(Z_A\subseteq A_+\) and \(Z_B\subseteq B_+\) be witnessing
families of minimum sizes. For every \((x,\rho)\), choose
\(a\in Z_A\) with \(A(x,a)\), and then \(b\in Z_B\) with
\(B(\rho(a),b)\). Thus \(Z_A\times Z_B\) witnesses the upper
product bound.

Conversely let \(T\subseteq A_+\times B_+\) witness the composition
norm. Put \(T_a=\{b:(a,b)\in T\}\) and
\(X=\{a:|T_a|\geq\|\mathbf B\|\}\). If \(|X|<\|\mathbf A\|\), choose
\(x\in A_-\) unrelated to every member of \(X\). For every
\(a\notin X\), the family \(T_a\) is too small to witness the
\(\mathbf B\)-norm; choose \(y_a\in B_-\) such that
\(\neg B(y_a,b)\) for all \(b\in T_a\). Define \(\rho(a)=y_a\)
there, and give \(\rho\) any values on \(X\). A witness
\((a,b)\in T\) for \((x,\rho)\) must have \(a\notin X\), because
\(A(x,a)\). But then \(b\in T_a\) contradicts
\(B(\rho(a),b)\). Hence \(|X|\geq\|\mathbf A\|\), and each member of
\(X\) has at least \(\|\mathbf B\|\) pairs in \(T\). Choice gives
\(|T|\geq\|\mathbf A\|\cdot\|\mathbf B\|\), including finite norms.

For the dual upper bound, let \(U\subseteq A_-\) witness
\(\|\mathbf A^\perp\|\) and fix any \(\rho_0:A_+\to B_-\). The family
\(\{(x,\rho_0):x\in U\}\) is unbounded for the composition:
against any \((a,b)\), one of its members has \(\neg A(x,a)\).
Alternatively let \(V\subseteq B_-\) witness \(\|\mathbf B^\perp\|\),
fix \(x_0\in A_-\), and for \(y\in V\) let \(\rho_y\) be the
constant function with value \(y\). Against \((a,b)\), choose
\(y\in V\) with \(\neg B(y,b)\). Then
\((x_0,\rho_y)\) is not related to \((a,b)\). These two families
give the minimum upper bound.

Finally suppose \(U\subseteq A_-\times{}^{A_+}B_-\) has size below
both dual norms. Its first-coordinate projection has a common
\(A\)-answer \(a\). The set \(\{\rho(a):(x,\rho)\in U\}\)
has a common \(B\)-answer \(b\). Thus \((a,b)\) answers every
member of \(U\); \(U\) cannot be a dual witness. This proves the
opposite minimum inequality and completes both equalities.
\end{proof}
\end{annexitem}

\begin{annexitem}{annex:p06.ccsep.domain-bridge}{thm:p06-ccsep-domain-bridge}{Measure-class and category-class transfers}
\begin{proof}
We first justify the measure transfer for arbitrary subsets. The binary
expansion \(\pi:2^\omega\to[0,1]\) is continuous and onto. A
cylinder \([s]\) has image a closed dyadic interval of length
\(2^{-|s|}\). Each dyadic point has at most two expansions, each a
singleton of product measure zero. The preimage of a half-open dyadic
interval differs from its cylinder only at these endpoint expansions,
so its measure equals the interval length. Countable disjoint dyadic
interval approximations from within any open interval show that its
preimage has measure equal to its length within \([0,1]\).

If \(E\subseteq[0,1]\) is outer null, cover it by open intervals of
summed length less than \(\epsilon\). The preimages cover
\(\pi^{-1}E\) and have summed measures at most \(\epsilon\), so
\(\pi^{-1}E\) is outer null. Conversely every null
\(A\subseteq2^\omega\) has an open cover of measure less than
\(\epsilon\). An open subset of Cantor space is the disjoint union
of its minimal cylinders; their weights sum to the open set's measure.
The corresponding dyadic image intervals cover \(\pi[A]\) with
summed lengths less than \(\epsilon\). Thus images of outer-null
sets are outer null too. No measurability of \(A\) or \(E\) is
used.

This two-way null preservation transfers the four minima. A family of
null sets with nonnull union has images with nonnull union: otherwise
the preimage of that null union would contain the original union.
The reverse transfer uses preimages and surjectivity. Images and
preimages also transfer covers. A nonnull Cantor set has a nonnull
image of no larger size. For a nonnull set in \([0,1]\), choose one
expansion of every point; the chosen set is nonnull, since its image
would otherwise be null. These give the two uniformity comparisons.
For cofinality, images of a cofinal Cantor base dominate every interval
null set, because its null preimage lies in a base member. Preimages
of a cofinal interval base dominate every Cantor null set, by first
taking that set's null image. Thus all four values are equal.

Removing a small set from a space does not change these four invariants
for the ideals at issue. More explicitly, if \(E\) belongs to a proper
sigma ideal and \(Y=X\setminus E\), intersecting witnesses with
\(Y\) preserves non-small unions and non-small sets, while adjoining
\(E\) to covering or cofinal witnesses restores the whole space.
The possible extra one member does not alter the infinite cardinals;
all four are uncountable here. The reverse comparisons use inclusion
and intersection. For nullity, the relative ideal on \(Y\) is simply
the ambient ideal restricted to \(Y\).

The homeomorphism
\[
 h:\mathbb R\to(0,1),\qquad
 h(t)=\tfrac12+\frac{t}{2(1+|t|)}
\]
and its inverse preserve null sets. Indeed \(h\) is Lipschitz, and
its inverse is Lipschitz on every compact interval
\([1/k,1-1/k]\), \(k\geq3\). Covering a null set by arbitrarily
short intervals proves that a Lipschitz image is null; a countable
union of these inverse restrictions is still null. Removing \(0,1\)
from \([0,1]\) and applying this homeomorphism therefore transfers
the interval values to \(\mathbb R\).

For category remove the countable set \(E\) of eventually constant
binary sequences and the dyadic points \(D\) in \((0,1)\).
They are meager, and their complements are dense \(G_\delta\)
subspaces. The restriction of \(\pi\) is a homeomorphism from
\(2^\omega\setminus E\) onto \((0,1)\setminus D\).
Continuity is inherited from \(\pi\); its inverse is continuous
because the first \(n\) digits are constant in a sufficiently small
neighborhood of any nondyadic point.

For a dense subspace \(Y\) of a space \(X\), a relatively nowhere
dense subset \(F\subseteq Y\) has nowhere dense closure in \(X\):
an open part of that closure would intersect the dense \(Y\) in
an open part of \(\overline F^{\,Y}\). Conversely the intersection
with \(Y\) of an ambient closed nowhere dense set is relatively
nowhere dense. Hence relative and ambient meagerness agree on
subsets of \(Y\). The small-complement argument just given transfers
the four invariants across these two dense subspaces, and the
homeomorphism transfers their meager ideals. Finally \(h\) transfers
category from \((0,1)\) to \(\mathbb R\). This completes all
four category comparisons as well as all four measure comparisons.
\end{proof}
\end{annexitem}

\begin{annexitem}{annex:p06.ccsep.splitting-category}{thm:p06-ccsep-splitting-category}{The comeager set of splitters}
\begin{proof}
Fix infinite \(Y\subseteq\omega\). For each \(n\), the sets of
reals identically zero on \(Y\setminus n\), or identically one on
\(Y\setminus n\), are closed nowhere dense. They are closed by
their coordinate equations. Given any finite condition, choose an
unused coordinate in \(Y\setminus n\) and assign the opposite value
to see that no cylinder is contained in either set. A real that does
not split \(Y\) belongs to the union of these two kinds of sets
over \(n<\omega\); conversely membership in that union means
one side of the split is finite. The nonsplitters therefore form a
meager set. Any nonmeager family has a member outside that set,
proving that it splits every infinite \(Y\).

For a family \(\mathcal R\) of infinite sets of size less than
\(\operatorname{cov}(\mathcal B)\), the meager nonsplitter sets
indexed by its members cannot cover Cantor space. A real outside
their union splits every member of \(\mathcal R\). Such a family
cannot witness \(\mathfrak r\). This proves both comparisons.
\end{proof}
\end{annexitem}

\begin{annexitem}{annex:p06.ccsep.chopped-base}{thm:p06-ccsep-chopped-base}{Finite-prefix uniformization of dense open sets}
\begin{proof}
The matching set is
\(\bigcap_m\bigcup_{n\geq m}\{y:y\mathbin{\upharpoonright}I_n
=x\mathbin{\upharpoonright}I_n\}\). Each tail union is open
dense: a finite stem can be extended to agree on a later whole
block. Thus the intersection is a dense \(G_\delta\).

Write a given dense \(G_\delta\) as \(\bigcap_nG_n\) with the
\(G_n\) decreasing dense open. Set \(a_0=0\). Given \(a_n\),
enumerate the finitely many binary strings \(t\) of length \(a_n\).
Construct one finite word \(w\), initially empty, by visiting them
one at a time. At the visit to \(t\), density and openness let
us append a word \(v\) so that \([t\frown w\frown v]\subseteq G_n\).
Replace \(w\) by \(w\frown v\). Previously obtained inclusions
are retained under further extension. After all visits, append a
bit if needed so that \(w\) is nonempty. Put
\(a_{n+1}=a_n+|w|\) and assign \(x\) on this block to be \(w\).
Every real agreeing with \(x\) on this block belongs to \(G_n\),
irrespective of its earlier prefix. If it matches on infinitely
many blocks, then for every \(m\) it belongs to some
\(G_n\subseteq G_m\) with \(n\geq m\); it therefore belongs
to the original intersection. This gives the required matching
subset. Apply it to the comeager complement of a closed nowhere
dense countable cover of a meager set for the last assertion.
\end{proof}
\end{annexitem}

\begin{annexitem}{annex:p06.ccsep.engulfing}{thm:p06-ccsep-engulfing}{Containment and the separating real}
\begin{proof}
If \((x,\Pi)\) engulfs \((x',\Pi')\), infinitely many matching
\(\Pi\)-blocks for a real \(y\) yield infinitely many contained
\(\Pi'\)-blocks on which \(y=x=x'\). Distinct containing blocks
yield distinct contained blocks, so \(y\) matches \((x',\Pi')\).

If engulfing fails, infinitely many \(\Pi\)-blocks are bad: they
contain no \(\Pi'\)-block on which \(x=x'\). Thin these to an
infinite family \(\mathcal I\) such that no \(\Pi'\)-block
meets two chosen blocks. This thinning is possible recursively:
the finitely many \(\Pi'\)-blocks meeting a fixed finite
\(\Pi\)-block end at a finite natural number; choose the next
bad block beyond all those endpoints for the previously selected
blocks. Set \(y=x\) on \(\bigcup\mathcal I\), and
\(y(k)=1-x'(k)\) elsewhere. The real \(y\) matches \((x,\Pi)\)
on every selected block. A \(\Pi'\)-block disjoint from their
union has a coordinate of disagreement with \(x'\). One meeting
a selected block but not contained in it has such a coordinate
outside the union, because it meets at most one selected block.
If it is contained in a selected block, badness gives a coordinate
with \(x\ne x'\), and \(y=x\) there. Thus no \(\Pi'\)-block
matches at all, and containment fails. This proves the converse.
\end{proof}
\end{annexitem}

\begin{annexitem}{annex:p06.ccsep.category-norms}{thm:p06-ccsep-category-norms}{Reading the four witness properties}
\begin{proof}
The data satisfy the standing norm assumptions. A chopped real is
matched by its own real and engulfed by itself. No single real matches
every chopped real: its bitwise opposite, chopped into singleton blocks,
is a counterexample. No single chopped real engulfs every chopped real
either: its complementary matching set is meager, whereas containing
every complementary matching set would, by the basis lemma applied to
each singleton, require it to contain the whole Cantor space.
A set of reals witnessing \(\|\mathbf U\|\) meets every matching set.
It is nonmeager: if it were meager, the matching-basis lemma would
supply a matching set disjoint from it. Conversely a nonmeager
set meets every comeager matching set. This proves the first equality.
A dual \(\mathbf U\)-witness is a family of chopped reals with no common
matching real, equivalently a family of their meager complementary
sets covering Cantor space. Every meager cover can be enlarged to
one by these complementary sets, by the basis lemma, so its least
size is precisely \(\operatorname{cov}(\mathcal B)\).

Write \(E_{x,\Pi}=2^\omega\setminus\operatorname{Match}(x,\Pi)\).
The engulfing criterion says
\(V((x,\Pi),(z,\Theta))\) exactly when
\(E_{x,\Pi}\subseteq E_{z,\Theta}\). These sets form a cofinal
base for the meager ideal. A \(\mathbf V\)-witness is therefore exactly
a cofinal family from this base. An arbitrary cofinal meager family
can be enlarged to such a family without increasing its size.
This proves the cofinality equality. Similarly a family of chopped
reals has a common engulfing chopped real exactly when the union
of its complementary matching sets is meager: one direction uses
that engulfing set as container, and the other uses the basis lemma
to enlarge the meager union, followed by the containment criterion.
The minimum size with no common answer is consequently
\(\operatorname{add}(\mathcal B)\), proving the fourth equality.
\end{proof}
\end{annexitem}

\begin{annexitem}{annex:p06.ccsep.category-sequential}{thm:p06-ccsep-category-sequential}{The answer-dependent second task and its Borel coding}
\begin{proof}
Code a chopped real by its binary real and increasing sequence of
endpoints. These codes form a Borel subset of the product of the
two sequence spaces, since the endpoint conditions use only integer
equalities and inequalities. Matching is a Borel relation, expressed
by infinitely many finite-block equations. For a chopped real
\(z=(x,\Pi)\) and a real \(y\) that matches it, define
\(\beta(z,y)(n)\) to be the right endpoint of the first matching
\(\Pi\)-block strictly after the block containing \(n\). This exists
because infinitely many blocks match. If \(y\) does not match,
put \(\beta(z,y)\) equal to the constant zero function, making
the map total.

For \(g\in\omega^\omega\), define the partition \(\Theta_g\)
by \(b_0=0\),
\(b_{j+1}=1+\max\{b_j,g(b_j)\}\), and put
\(\gamma(y,g)=(y,\Theta_g)\). Define the negative morphism
component by
\[
 z\longmapsto(z,\rho_z),\qquad \rho_z(y)=\beta(z,y),
\]
and the positive component by \((y,g)\mapsto\gamma(y,g)\).
If \(U(z,y)\) and \(\beta(z,y)\leq^*g\), then for every
sufficiently late \(\Theta_g\)-block starting at \(b_j\), the
matching \(\Pi\)-block selected by \(\beta(z,y)(b_j)\) starts
at or after \(b_j\) and ends at most at
\(g(b_j)<b_{j+1}\). Thus the new block contains an old block
on which \(y=x\). The chopped real \(\gamma(y,g)\) engulfs
\(z\). This is exactly the morphism implication for \(\mathbf U;\mathbf W\to\mathbf V\).

The identity first component is continuous. Each coordinate of
\(\beta\), on the matching part of its domain, is obtained by
a least finite-block search; the condition that its value is a
particular endpoint is a finite-block condition and finitely many
preceding failures. The matching-domain test is Borel, and the
nonmatching branch is constant. Thus the uncurried \(\beta\) is
Borel. Every finite list of endpoints of \(\Theta_g\) depends
on only finitely many values of \(g\), by its finite recursion.
The endpoint map and \(\gamma\) are therefore continuous.
These are the required Borel component claims; no Borel structure
on the whole higher-type function space is asserted.

Apply the morphism-duality and sequential-norm results. The norm
comparison gives
\(\|\mathbf V\|\leq\|\mathbf U\|\cdot\|\mathbf W\|\), which is the displayed
maximum because both factors are infinite cardinals. The dual
comparison gives
\(\|\mathbf V^\perp\|\geq\min\{\|\mathbf U^\perp\|,\|\mathbf W^\perp\|\}\).
The category-norm result substitutes the four ideal cardinals,
and \(\|\mathbf W\|=\mathfrak d\), \(\|\mathbf W^\perp\|=\mathfrak b\)
substitute the function cardinals. The opposite bounds are already
proved in the classical inequality network. The domain bridge
therefore recovers the same numerical category identities while
retaining this separate Borel witness presentation.
\end{proof}
\end{annexitem}

\begin{annexitem}{annex:p06.ccsep.support}{thm:p06-ccsep-support}{Antichains, nice names, and product factorization}
\begin{proof}
First recall the finite-set delta-system argument, in the precise
uncountable form needed here. Thin an uncountable family of finite
sets to a fixed size \(n\). Induct on \(n\). If some element
belongs to uncountably many sets, restrict to those sets, remove it,
and use the induction hypothesis, then restore it to the root.
If every element belongs to only countably many sets, recursively
choose \(\omega_1\) pairwise disjoint members: at each countable
stage the previously used elements forbid only countably many family
members. This is a delta system with empty root. The induction
also includes size zero.

Apply this to the finite domains of an uncountable collection of
Cohen conditions. Their common root has only finitely many binary
assignments, so further thin to conditions with the same root values.
Every two now have a consistent union, a common strengthening. There
is no uncountable antichain. Cardinal and cofinality preservation
follow from the already proved ccc preservation theorem
\cref{thm:p06-ccc-preservation}.

For a name for a real, for each natural-number coordinate take a
maximal antichain deciding its value (or its bit in a binary code).
Each antichain is countable. The union of the finite coordinate
supports of all these conditions is a countable ground set \(S\).
The decisions, regarded as a nice name on \(P_S\), determine the
same real. A countable sequence of codes is handled by countably
many such antichains; the union is still countable. For a family
enumerated by a ground ordinal \(\lambda<\kappa\), choose these
names for every index and take their support union. Its ground
size is at most \(|\lambda|\cdot\aleph_0<\kappa\). If the
family is first specified in the extension, apply the forcing
maximum principle to a name for its enumeration below a condition
that forces its length. Cardinal preservation gives a ground
ordinal \(\lambda<\kappa\) for this length; the selected names and
their union are then ground objects. Thus this argument does not
mistake an extension family of supports for a ground set.

The restriction map sends \(P_\kappa\) isomorphically to
\(P_S\times P_{\kappa\setminus S}\): split each finite function
into its two restrictions, and combine by union for the inverse.
The projected filter is generic in the first factor. The second
projection is generic over \(M[G_S]\). To verify the latter,
given a first-factor name forced to be a dense subset of the second
factor, extend any first component to decide a second condition
below its present second component that belongs to that named dense
set; the deciding pairs are dense in the product. The generic meets
them. If density is forced only below a first-factor condition in
the generic, perform the same argument below that condition; add
conditions incompatible with it to obtain a full dense set.
Projection of the second factor to one unused coordinate gives a
Cohen real over \(M[G_S]\) by the same product argument.

There are at most \((\kappa^{\aleph_0})^{\aleph_0}=\kappa\)
binary nice names: a name is given by countably many countable
antichains in the poset of ground size \(\kappa\), with bit labels.
Thus the extension has at most \(\kappa\) reals. For distinct
\(\alpha,\beta<\kappa\), conditions making their bits differ at
some previously unused natural coordinate are dense. Therefore all
\(c_\alpha\) are distinct, giving at least \(\kappa\) reals.
This proves \(\mathfrak c=\kappa\).
\end{proof}
\end{annexitem}

\begin{annexitem}{annex:p06.ccsep.category-preservation}{thm:p06-ccsep-category-preservation}{A ground comeager set tailored to the new cover}
\begin{proof}
Suppose a condition \(p\) in a countable forcing forces that
\(A\) is covered by closed nowhere dense sets \(\dot F_n\),
\(n<\omega\). For each \(q\leq p\) and \(n\), define in
the ground
\[
 H_{q,n}=\bigcup\{[s]:\exists r\leq q\quad
                   r\forces [s]\cap\dot F_n=\varnothing\}.
\]
This set is open dense. Indeed, given a cylinder \([t]\), \(q\)
forces that the nowhere dense \(\dot F_n\) misses a subcylinder
\([s]\subseteq[t]\). Some strengthening decides such a finite
\(s\) and forces the required disjointness. Thus \([t]\) meets
\(H_{q,n}\).

There are only countably many pairs \((q,n)\). Since \(A\) is
nonmeager in the ground, choose
\(x\in A\cap\bigcap_{q\leq p,n<\omega}H_{q,n}\) there.
For each \(n\) and each \(q\leq p\), a strengthening \(r\)
forces \(x\notin\dot F_n\), by the definition of \(H_{q,n}\)
and the chosen membership of \(x\). These strengthenings are
dense below \(p\). The density property of forcing therefore gives
\(p\forces x\notin\dot F_n\), for each \(n\), contradicting
the forced cover. Any meager set has a closed nowhere dense
countable cover; hence no new meager set can contain \(A\).

In the full Cohen extension a proposed meager container of
\((2^\omega)^M\) has a countable sequence of closed nowhere dense
codes. By the support lemma these codes lie in \(M[G_S]\) for
a countable ground \(S\). The poset \(P_S\) is countable in
\(M\). The coded sets are still closed nowhere dense in this
intermediate model. For a closed set coded by a binary tree, being
nowhere dense says that every finite cylinder has a subcylinder
disjoint from the tree; this property uses the same finite strings
and integers in both extensions and is absolute. The countable
forcing argument thus gives a ground real outside the proposed
cover in \(M[G_S]\). Its nonmembership in each coded closed set
is also absolute, so it stays outside in \(M[G]\).

The old Cantor space is nonmeager in \(M\) by Baire's theorem and
has size \(\aleph_1\) by ground CH. Cardinal preservation retains
that size. It is therefore a witness to
\(\operatorname{non}(\mathcal B)\leq\aleph_1\). The proper
sigma-ideal lower bound gives equality.
\end{proof}
\end{annexitem}

\begin{annexitem}{annex:p06.ccsep.fresh-obstructions}{thm:p06-ccsep-fresh-obstructions}{Avoiding, escaping domination, and splitting}
\begin{proof}
For a family of fewer than \(\kappa\) meager sets, choose a
closed nowhere dense cover of each one. All these codes lie in
\(M[G_S]\) with \(|S|^M<\kappa\) by the support lemma. Choose
\(\alpha\notin S\). The Cohen real \(c_\alpha\) avoids every
one of these coded meager sets. For any closed nowhere dense set
\(F\) in that model, conditions whose cylinder misses \(F\)
are dense. Meeting these dense sets for all closed pieces puts
\(c_\alpha\) outside the proposed cover. The correctness of the
codes is absolute as in the preceding proof. Hence
\(\operatorname{cov}(\mathcal B)\geq\kappa\).

For domination take \(D\subseteq\omega^\omega\) of size less
than \(\kappa\), and put it inside such a smaller extension.
A fresh \(c=c_\alpha\) has infinitely many ones: conditions
putting a one beyond any specified coordinate are dense. Let
\(f_c(n)\) be the position of its \(n\)-th one, counted from
zero. Fix \(g\in\omega^\omega\) in the smaller extension and
an integer \(N\). The requirement
\(\exists n\geq N\ f_c(n)>g(n)\) is dense. To see this,
extend a finite condition to an initial binary stem with \(k\)
ones. Choose \(n\geq\max\{N,k\}\), add ones until the stem
has exactly \(n\) ones, then add zeros until its next position
exceeds \(g(n)\), and put a one there. This fixes the position
of the \(n\)-th one above \(g(n)\). Genericity for every \(N\)
gives \(f_c\not\leq^*g\). Since this holds for every \(g\in D\),
the family \(D\) does not dominate \(f_c\). Thus
\(\mathfrak d\geq\kappa\).

For reaping take a family \(\mathcal R\subseteq[\omega]^\omega\)
of size below \(\kappa\) in the smaller extension. For each
\(Y\in\mathcal R\), each integer \(N\), and each bit
\(i\in\{0,1\}\), conditions assigning value \(i\) to some
unused \(n\in Y\setminus N\) are dense. Infinitude of \(Y\)
is absolute. The fresh coordinate meets all these ground dense
sets, so its set of ones splits every member of \(\mathcal R\).
No such family is reaping, giving \(\mathfrak r\geq\kappa\).

All three invariants are at most the continuum, by singleton covers
for category, all functions for domination, and all infinite subsets
for reaping. The support lemma computed that continuum as
\(\kappa\), completing the three equalities.
\end{proof}
\end{annexitem}

\begin{annexitem}{annex:p06.ccsep.null-values}{thm:p06-ccsep-null-values}{A summable cover from block matches}
\begin{proof}
For fixed \(c\), the measure of the union of its matching-block
events from index \(m\) onward is at most
\(\sum_{n\geq m}2^{-n-1}=2^{-m}\). Their intersection over
\(m\) has outer measure zero. Thus \(N_c\) is null in any
model containing \(c\), by the same explicit bound.

Let \(c\) be Cohen over \(M'\), and fix \(x\in(2^\omega)^{M'}\).
For every \(m\), conditions agreeing with \(x\) on some whole
block \(I_n\), \(n\geq m\), are dense: choose such a block
disjoint from the finite condition domain, and assign its finitely
many bits from \(x\). Meeting these dense requirements gives
infinitely many matching blocks, so \(x\in N_c\). This holds
for every ground \(x\), not merely for countably many of them.

If \(X\subseteq2^\omega\) has size less than \(\kappa\) in the
full extension, it is contained in some \(M[G_S]\) with
\(|S|^M<\kappa\). A fresh \(c_\alpha\) is Cohen over this
model, and \(X\subseteq N_{c_\alpha}\). The set is null by
the uniform tail bound, so \(X\) cannot witness uniformity.
This proves \(\operatorname{non}(\mathcal N)\geq\kappa\);
the whole Cantor space supplies the opposite bound.

For covering use just the family
\(\{N_{c_\alpha}:\alpha<\omega_1^M\}\). Given any real \(x\)
in the full extension, choose a countable ground support \(S\)
for it. Some \(\alpha<\omega_1^M\) is outside \(S\). The
coordinate \(c_\alpha\) is Cohen over the smaller model
containing \(x\); therefore \(x\in N_{c_\alpha}\). These
\(\aleph_1\) null sets cover Cantor space. No countable null
family can cover a probability-one space, since its union remains
null. Hence \(\operatorname{cov}(\mathcal N)=\aleph_1\).
\end{proof}
\end{annexitem}

\begin{annexitem}{annex:p06.ccsep.cohen-values}{thm:p06-ccsep-cohen-values}{Computing every displayed value rather than just the continuum}
\begin{proof}
The support, preservation, obstruction, and null-value lemmas give
\[
 \mathfrak c=\mathfrak d=\mathfrak r
 =\operatorname{cov}(\mathcal B)
 =\operatorname{non}(\mathcal N)=\kappa,
 \qquad
 \operatorname{non}(\mathcal B)
 =\operatorname{cov}(\mathcal N)=\aleph_1.
\]
The domain bridge carries the four ideal values to the real-line
domains of the classical inequality network. The remaining entries
are then forced by proved bounds, as follows. The inequalities
\(\aleph_1\leq\mathfrak b\leq\operatorname{non}(\mathcal M)\)
give \(\mathfrak b=\aleph_1\). The splitting-category result and
the countable splitting obstruction give \(\mathfrak s=\aleph_1\).
Since \(\operatorname{add}(\mathcal M)\leq\mathfrak b\) and
the meager ideal is a proper sigma ideal, its additivity is
\(\aleph_1\). Likewise
\(\aleph_1\leq\operatorname{add}(\mathcal N)
\leq\operatorname{cov}(\mathcal N)=\aleph_1\).
For cofinality use
\(\mathfrak d\leq\operatorname{cof}(\mathcal M)\leq\mathfrak c\)
and
\(\operatorname{non}(\mathcal N)\leq\operatorname{cof}(\mathcal N)
\leq\mathfrak c\), obtaining \(\kappa\) in both cases.
These bounds are supplied by
\cref{thm:p06-cc-combinatorial,thm:p06-cc-ideal-bounds,thm:p06-cc-cichon},
not inferred from a diagram drawing without proofs.

For the stated \(\mathsf{GCH}\) example take
\(\kappa=\aleph_2\). Successor regularity in this case follows
from \(\aleph_2=2^{\aleph_1}\) and Koenig's theorem:
\(\operatorname{cf}(2^{\aleph_1})>\aleph_1\), while that
cofinality is at most \(\aleph_2\). Furthermore
\[
 (\aleph_2)^{\aleph_0}
  =(2^{\aleph_1})^{\aleph_0}
  =2^{\aleph_1\cdot\aleph_0}
  =2^{\aleph_1}=\aleph_2.
\]
Thus the standing cardinal hypotheses are satisfied. Preservation
retains the distinction \(\aleph_1<\aleph_2\) in the extension,
and every listed strict comparison follows from its two explicitly
computed values.
\end{proof}
\end{annexitem}
